\documentclass[11pt]{article}

\usepackage[margin=1in]{geometry}
\usepackage{amsmath,amssymb,amsthm}
\usepackage{bm}
\usepackage{booktabs}
\usepackage{enumitem}
\usepackage[hidelinks]{hyperref}

\newtheorem{remark}{Remark}
\newtheorem{proposition}{Proposition}

\DeclareMathOperator*{\argmax}{arg\,max}
\DeclareMathOperator{\E}{\mathbb{E}}
\DeclareMathOperator{\Var}{Var}
\DeclareMathOperator{\Cov}{Cov}
\DeclareMathOperator{\tr}{tr}
\newcommand{\R}{\mathbb{R}}
\newcommand{\Norm}{\mathcal{N}}
\newcommand{\T}{^{\mathsf{T}}}
\newcommand{\given}{\,|\,}
\newcommand{\dd}{\mathrm{d}}
\newcommand{\half}{\tfrac{1}{2}}

\title{On the fly calibration and imaging of streamed TART data}
\author{Landman Bester}
\date{\today}

\begin{document}
\maketitle

\begin{abstract}
This note develops the calibration and imaging machinery underpinning the proposed transient detector software pipeline; the detection heuristics themselves are derived in an accompanying document. Because those heuristics apply to calibrated, source-subtracted residual images, on-the-fly direction-independent calibration and calibrator subtraction are a necessary precursor. TART is a single-polarisation, zenith-pointing all-sky array, so we track one scalar complex gain per antenna, written as $g_p=\exp(a_p+\imath\psi_p)$, together with one real, non-negative flux per catalogued calibrator. The slowly varying log-amplitudes $a_p$, the faster varying phases $\psi_p$, and the source fluxes $A_s$ are each modelled as independent integrated Wiener processes (IWPs); see the accompanying document for IWP notation and theory. The driving variances are \emph{not} inferred online---they are calibrated once, offline, from the large body of existing data reductions---which avoids the ill-conditioned coupling between hyperparameter inference and a heavy-tailed likelihood. Conditioned on the gains the model visibility is \emph{linear} in the fluxes, so the fluxes are Rao--Blackwellised out in closed form while the gains are updated by a robust iterated extended Kalman filter: an antenna-based Gauss--Newton step in $(a,\psi)$, \emph{not} a complex-gain StEFCal iteration. Robustness to the inevitably incomplete satellite catalogue comes from a heavy-tailed (Student-$t$) observation model, realised as iteratively reweighted least squares. The filter emits gain-corrected, model-subtracted residual visibilities, which are imaged onto a HEALPix grid to produce the residual image stream consumed by the detector.
\end{abstract}

\section{Overview of the imaging and calibration problem}

The task is to transform the raw uncalibrated visibilities $V_{pq,t}(\bm{b})$ into calibrated residual images $I_{t}(\bm{l})$ in real time, where $p,q$ index the antennas, $t$ the time, $\bm{b}$ the baseline vector, and $\bm{l}$ the direction on the sky. Because real-time streaming is required, calibration and calibrator source flux estimation needs to happen simultaneously (it is only the calibrator positions which are known a-priori and these are not always 100\% accurate). The streaming nature of the problem makes Holoscan a natural fit for the software implementation, with CuPy---rather than an autodiff framework such as JAX---as the numerical backend for the inference machinery. We prefer CuPy because the pipeline must eventually run on constrained hardware, where a single, explicitly budgeted GPU memory pool is important, and because the calibration Jacobian is highly structured (Section~\ref{sec:obs}), so automatic differentiation would buy little (Section~\ref{sec:notes}).

We assume a measurement model of the form
\begin{equation}
  V_{pq,t}(\bm{b}) = \exp\left(a_p(t) + \imath \psi_p(t)\right) \left( \sum_s A_s(t) X_{pq,ts} \right) \exp\left(a_q(t) - \imath \psi_q(t)\right),
\end{equation}
where the scalar gain log-amplitudes $a_{p/q}$ are slowly varying real functions of time, the gain phases $\psi_{p/q}$ are more quickly varying real functions of time, the calibrator source fluxes $A_s$ are slowly varying real functions of time, and $X_{pq,ts}$ is the known unit coherencies of source $s$ to the visibility on baseline $pq$ at time $t$. The TART telescope~\cite{tart} always points at zenith and sees almost the whole hemisphere it is located in. This makes spherical coordinates the natural choice for the image domain and the calibrator source positions. The residual images should be computed on a Healpix grid.

Note that calibrator sources are mostly satellites in geostationary orbit, which are bright and have well-known positions but whose fluxes can vary on time-scales of minutes to hours. The gain log-amplitudes and phases are driven by independent integrated Wiener processes (IWPs) with different driving noise levels. The calibrator source fluxes are also driven by independent IWPs, but with much smaller driving noise levels than the gain parameters. Since the calibrator satellite catalogue is almost always incomplete, a robust heavy-tailed likelihood is required to prevent outliers from dominating inference. The IWP state-space model and the heavy-tailed likelihood are combined in a robust iterated extended Kalman filter: the gains are updated by an antenna-based Gauss--Newton step in $(a,\psi)$, while the catalogue fluxes---which enter the model linearly---are Rao--Blackwellised out in closed form. The filter tracks the calibration parameters, subtracts the calibrator sources, and emits a gain-corrected residual visibility stream that is imaged to produce the residual images fed to the detection pipeline.

The remainder of this note formalises the calibration filter and the imaging step. Section~\ref{sec:statespace} sets up the IWP state-space model and fixes the gauge; Section~\ref{sec:obs} gives the linearised (iterated-EKF) observation update and its Jacobian; Section~\ref{sec:robust} adds the heavy-tailed likelihood; Section~\ref{sec:flux} Rao--Blackwellises the source fluxes and describes their warm-up; Section~\ref{sec:imaging} defines the gain-corrected residual visibilities and their projection onto a HEALPix grid; Section~\ref{sec:hyper} covers offline calibration of the driving variances; and Section~\ref{sec:notes} collects implementation choices and reuse of existing TART software (the \texttt{tart-telescope} organisation and T.~Molteno's repositories, especially \textsc{disko} and \textsc{spotless}).

\section{State-space model for the calibration parameters}
\label{sec:statespace}

Let $P$ be the number of antennas and $S_t$ the number of catalogued sources above the
horizon at time $t$. We track the gain log-amplitudes $a_p$, the gain phases $\psi_p$
($p=1,\dots,P$), and the source fluxes $A_s$ ($s=1,\dots,S_t$). Each scalar parameter is an
independent IWP~\cite{sarkka}; following the accompanying note we lift a scalar parameter $\theta$ and
its first $q$ derivatives into a state $\bm{\theta}=(\theta,\dot\theta,\dots,\theta^{(q)})\T$
with the exact discrete transition
\begin{equation}
  \bm{\theta}_k = \bm{A}_\theta(\Delta_k)\,\bm{\theta}_{k-1}+\bm{w}_k,
  \qquad \bm{w}_k\sim\Norm(\bm 0,\bm{Q}_\theta(\Delta_k)),
\end{equation}
where $\bm{A}_\theta,\bm{Q}_\theta$ are the closed-form IWP matrices of the accompanying
document and depend on the per-parameter driving variance $\sigma_\theta^2$. We default to
$q=1$ (a locally linear, ``constant-velocity'' model) for every parameter; the order may be
raised per group if the offline data warrant it. Stacking the per-parameter blocks gives the
full state $\bm{x}_k$ and a block-diagonal transition
\begin{equation}
  \bm{x}_k=\bm{F}_k\,\bm{x}_{k-1}+\bm{w}_k,
  \qquad \bm{F}_k=\operatorname{blkdiag}\big(\{\bm{A}_{a_p}\},\{\bm{A}_{\psi_p}\},\{\bm{A}_{A_s}\}\big),
\end{equation}
with $\bm{Q}_k$ the matching block-diagonal process covariance, so the predict step is exact,
decoupled across parameters, and $O(P+S_t)$ per frame.

\paragraph{Driving variances are fixed.} The driving-variance groups---$\sigma_a^2$ (slow),
$\sigma_\psi^2$ (faster), and $\sigma_A^2$ (slow, modulo the warm-up of
Section~\ref{sec:flux})---are \emph{not} estimated online. They are calibrated once from the
existing offline reductions (Section~\ref{sec:hyper}). This is deliberate: inferring the
driving variances through the marginal likelihood while a heavy-tailed likelihood
simultaneously reweights the data is poorly conditioned, the two mechanisms competing to
explain the same residual.

\paragraph{Identifiability and gauge fixing.} The likelihood depends on the gains only through
the products $g_p g_q^{\ast}=\exp(a_p+a_q)\exp\!\big(\imath(\psi_p-\psi_q)\big)$, which leaves
two exact gauge freedoms:
\begin{itemize}[leftmargin=*]
  \item a global phase $\psi_p\mapsto\psi_p+\phi$ cancels in every $g_pg_q^{\ast}$; we remove
        it by pinning a \emph{reference-antenna} phase, $\psi_{\mathrm{ref}}\equiv0$;
  \item a global log-amplitude $a_p\mapsto a_p+c$ rescales every visibility by $e^{2c}$ and is
        exactly degenerate with an overall flux rescaling $A_s\mapsto e^{-2c}A_s$. Because the
        catalogue supplies positions but \emph{no} fluxes, this degeneracy cannot be anchored
        on the source side, so we impose the zero-sum constraint $\sum_p a_p=0$.
\end{itemize}
Both constraints are enforced inside the update (Section~\ref{sec:obs}): the reference phase is
eliminated from the state, and the gain mean is projected onto $\sum_p a_p=0$ after each
Gauss--Newton step.

\section{Observation model and its linearisation}
\label{sec:obs}

At time $t_k$ each baseline $pq$ contributes a complex visibility. With the
single-polarisation scalar gains and the catalogue coherencies $X_{pq,s}$---unit-amplitude
fringes computed from the known source directions (Section~\ref{sec:imaging})---the noise-free
model is
\begin{equation}
  \hat V_{pq} = g_p\,\hat M_{pq}\,g_q^{\ast},\qquad
  \hat M_{pq}=\sum_{s} A_s\,X_{pq,s},\qquad
  g_p=\exp(a_p+\imath\psi_p).
  \label{eq:model}
\end{equation}
We observe $V_{pq}=\hat V_{pq}+n_{pq}$ with circular complex noise $n_{pq}$ whose variance is
set by the data weights, stacking real and imaginary parts as two real observations per
baseline.

\paragraph{Why not StEFCal.} StEFCal~\cite{stefcal} exploits that, holding all \emph{other} antennas' complex
gains fixed, $\hat V_{pq}$ is \emph{linear} in $g_p$, giving an antenna-by-antenna linear least
squares. Our state is $(a_p,\psi_p)$, not $g_p$, and $\hat V_{pq}$ is nonlinear in these, so we
linearise directly (Gauss--Newton / iterated EKF) rather than iterating a complex-gain fixed
point. The relevant structure nonetheless survives, and in a cleaner form.

\paragraph{Jacobian.} Differentiating \eqref{eq:model} at the current estimate,
\begin{equation}
  \frac{\partial \hat V_{pq}}{\partial a_p}=\frac{\partial \hat V_{pq}}{\partial a_q}=\hat V_{pq},
  \qquad
  \frac{\partial \hat V_{pq}}{\partial \psi_p}=+\imath\,\hat V_{pq},
  \qquad
  \frac{\partial \hat V_{pq}}{\partial \psi_q}=-\imath\,\hat V_{pq},
  \qquad
  \frac{\partial \hat V_{pq}}{\partial A_s}=g_p\,X_{pq,s}\,g_q^{\ast}.
  \label{eq:jac}
\end{equation}
The gain derivatives are simply the model visibility rotated by $0$ or $\tfrac{\pi}{2}$, so the
Jacobian costs nothing beyond $\hat V_{pq}$, which we already form. Each baseline row touches
only its two antennas: the gain block of the Jacobian inherits the \emph{bipartite}
(baseline$\,\times\,$antenna) sparsity of the array's incidence structure, whereas the flux
block $\partial\hat V_{pq}/\partial A_s$ is dense (every baseline sees every source). The
antenna-based normal-equations structure that StEFCal exploits therefore re-appears in
$\bm{H}_k^{\ast}\bm{H}_k$, despite the change of parametrisation---which is what makes a
hand-coded, structure-aware solver attractive.

\paragraph{Iterated-EKF measurement update.} Let $\bm{H}_k$ be the stacked (real/imag) Jacobian
\eqref{eq:jac} and $\bm{e}_k=V_k-\hat V_k(\bm{x}_{k|k-1})$ the innovation at the predicted mean.
One Gauss--Newton step is the standard EKF update~\cite{barshalom}; relinearising about the running estimate and
repeating yields the \emph{iterated} EKF, i.e. a Gauss--Newton solver for the MAP state under
the IWP prior. A few iterations suffice because the nonlinearity is mild. The normalised
innovation (NIS) formed here is logged as the whiteness diagnostic~\cite{kailath} of the accompanying note,
but---per Section~\ref{sec:statespace}---is never fed back to adapt the driving variances.

\section{Robust update via a heavy-tailed likelihood}
\label{sec:robust}

The Gaussian observation model of Section~\ref{sec:obs} is fragile: a baseline contaminated by an
un-catalogued source or by radio-frequency interference (RFI) is a gross outlier that, under a
Gaussian likelihood, would bias every gain and flux. We replace the Gaussian by a Student-$t$~\cite{studentt}
observation density, which has the Gaussian scale-mixture representation
\begin{equation}
  t_\nu(\bm{e}\given \bm 0,\bm{R}) = \int_0^\infty \Norm(\bm{e}\given\bm 0,\bm{R}/\lambda)\,
  \mathrm{Gamma}\!\big(\lambda\given\tfrac\nu2,\tfrac\nu2\big)\,\dd\lambda,
\end{equation}
in which the latent per-baseline precision scale $\lambda_{pq}$ down-weights discrepant data.
Treating $\{\lambda_{pq}\}$ by expectation--maximisation (equivalently, one variational step) gives
an iteratively reweighted least squares (IRLS) update: at Gauss--Newton iteration $j$ each baseline
carries a weight
\begin{equation}
  w_{pq}^{(j)} = \frac{\nu+n_y}{\nu + \epsilon_{pq}^{(j)}},
  \qquad
  \epsilon_{pq}^{(j)} = \bm{e}_{pq}^{(j)\,\T}\,\bm{R}_{pq}^{-1}\,\bm{e}_{pq}^{(j)},
\end{equation}
with $n_y=2$ for the stacked real/imag residual, and the EKF normal equations are assembled with the
reweighted observation precision $w_{pq}\bm{R}_{pq}^{-1}$. Inliers keep $w_{pq}\!\approx\!1$;
outliers are smoothly suppressed. The degrees-of-freedom $\nu$ is fixed (a small value, $\nu\sim4$,
is robust) rather than inferred, again to avoid the coupling of Section~\ref{sec:statespace}.

\begin{remark}[Two failure modes of an incomplete catalogue]
Robustness here and flux handling in Section~\ref{sec:flux} address \emph{different} symptoms of
catalogue incompleteness, and both are needed:
\begin{itemize}[leftmargin=*]
  \item \textbf{Present but un-catalogued} sources have no position, hence no $X_{pq,s}$, and cannot
        be modelled; they appear as heavy-tailed outliers and are handled \emph{only} by the
        Student-$t$ down-weighting.
  \item \textbf{Catalogued but absent} sources do have an $X_{pq,s}$ but no real emission; they are
        handled by driving their flux state to zero (Section~\ref{sec:flux}).
\end{itemize}
\end{remark}

\section{Rao--Blackwellised flux update and warm-up}
\label{sec:flux}

Conditioned on the gains, the model \eqref{eq:model} is \emph{linear} in the fluxes: with $g_p,g_q$
fixed, $\hat V_{pq}=\sum_s A_s\,(g_p X_{pq,s} g_q^{\ast})$. Partition the state into the nonlinear
gains $\bm{x}_g=\{a_p,\psi_p\}$ and the linear fluxes $\bm{x}_A=\{A_s\}$. The joint predicted density
is Gaussian (block-diagonal, the IWPs being independent) and the flux-conditional posterior
$p(\bm{x}_A\given\bm{x}_g,V_k)$ is exactly Gaussian. We therefore \emph{Rao--Blackwellise} the
fluxes: rather than sampling or jointly Newton-stepping them, we eliminate $\bm{x}_A$ in closed form
via the Schur complement of the (reweighted) normal-equations matrix, so the gain update sees the
fluxes marginalised with their full uncertainty and the flux posterior is recovered by
back-substitution. This lowers the variance of the gain estimate and confines the per-frame
nonlinear solve to the $2P$-dimensional gain block.

\paragraph{Non-negativity.} A physical flux satisfies $A_s\ge0$. We keep $A_s$ \emph{linear} (not
log-parametrised like the gains) so the Rao--Blackwell step stays exact, and impose non-negativity
as a constraint on the flux block, turning its update into a small non-negative least squares (a
truncated-Gaussian posterior); the Schur elimination then runs on the \emph{active} (strictly
positive) set each frame. Non-negativity is what makes a genuinely absent source harmless: lacking
any positive likelihood gradient, and forbidden from going negative to absorb noise, its flux
estimate simply stays near zero.

\paragraph{Warm-up schedule.} At acquisition the fluxes are unknown and the gains are not yet
locked, so we initialise every flux at mean zero with a diffuse covariance and run the flux driving
variance $\sigma_A^2$ \emph{large}, letting the data raise the fluxes of sources that are truly
present. As the filter converges---monitored via the settling of the gain-block NIS rather than a
fixed frame count---$\sigma_A^2$ is annealed down to the realistic floor calibrated offline
(Section~\ref{sec:hyper}); an exponential schedule is the fallback. Because a large warm-up
$\sigma_A^2$ also lets noise excite \emph{absent} sources (which then freeze when the variance
tightens), the zero-mean initialisation, the non-negativity constraint, and the Student-$t$ weights
together bias those fluxes back down. Should residual leakage into absent sources be observed, a
mild $\ell_1$ (LASSO-style proximal) shrinkage on $\bm{x}_A$ supplies \emph{active} suppression; we
start without it and add it only if needed.

\section{Calibrated residuals and all-sky imaging}
\label{sec:imaging}

\paragraph{Coherencies from the catalogue.} TART points at zenith and the satellite catalogue gives
positions (from broadcast ephemerides / two-line elements), not fluxes. For source $s$ with
time-dependent direction cosines $\bm{\hat l}_s(t)=(l_s,m_s,n_s)$ in the array frame and baseline
$\bm{b}_{pq}=(u_{pq},v_{pq},w_{pq})$, the unit coherency is the geometric fringe
\begin{equation}
  X_{pq,s}(t)=\exp\!\Big(2\pi\imath\,\tfrac{\nu}{c}\,
  \big[u_{pq}l_s+v_{pq}m_s+w_{pq}(n_s-1)\big]\Big),
  \qquad n_s=\sqrt{1-l_s^2-m_s^2}.
\end{equation}
Satellites in medium Earth orbit move appreciably over minutes, so $X_{pq,s}$ is recomputed every
frame from the propagated positions; the motion is in fact helpful, modulating each source's fringe
and so aiding the flux/gain separation. The coordinate pipeline (ephemeris propagation,
ECEF$\to$topocentric$\to$direction cosines, horizon masking) is the natural place to reuse TART
community code (Section~\ref{sec:notes}).

\paragraph{Gain-corrected residual visibilities.} After the update the filter forms the model
$\hat M_{pq}=\sum_s\hat A_s X_{pq,s}$ and emits, per baseline, the gain-corrected,
model-subtracted residual together with its corrected weight,
\begin{equation}
  R_{pq}=\hat g_p^{-1}\,V_{pq}\,\hat g_q^{-\ast}-\hat M_{pq},
  \qquad
  w^{\mathrm{corr}}_{pq}=w_{pq}\,\lvert \hat g_p\hat g_q\rvert^{2}.
  \label{eq:resid}
\end{equation}
Working in the gain-corrected (sky) frame puts the residual in physical flux units with a
well-characterised noise variance---exactly the per-frame variance $R_k$ the detector consumes. By
construction $R_{pq}$ has the same layout as a raw visibility, so the downstream imager need not know
that calibration occurred.

\paragraph{HEALPix imaging.} The residual map is the adjoint (gridless) DFT of $R_{pq}$ evaluated at
the centres of the visible HEALPix pixels~\cite{healpix} rather than on a tangent-plane $(l,m)$ grid:
\begin{equation}
  I_t(\bm{\hat l}_j)=\frac{1}{n_j}\,\mathrm{Re}\!\sum_{pq}
  w^{\mathrm{corr}}_{pq}\,R_{pq}\,
  \exp\!\Big(\!-2\pi\imath\,\tfrac{\nu}{c}\,[u_{pq}l_j+v_{pq}m_j+w_{pq}(n_j-1)]\Big),
\end{equation}
the natural representation for a zenith-pointing instrument that sees most of the hemisphere (a
single tangent plane breaks down at large zenith angle). The existing tiled $(l,m)$ DFT operator is a
structural template: the per-pixel phase, the $w$-term $(n_j-1)$ correction, and the memory tiling
all carry over once the pixel coordinates are HEALPix direction cosines. Each frame yields one
all-sky residual map; the stream of maps is the per-pixel light-curve input to the detector, whose
quiescent per-pixel noise level follows from $\{w^{\mathrm{corr}}_{pq}\}$.

\section{Offline calibration of the driving variances}
\label{sec:hyper}

The driving variances are estimated once, offline, from the existing body of data reductions and then
frozen for the online filter. Given an offline gain solution sampled at $\{t_k\}$, the IWP increments
of each log-amplitude and phase track are (for $q=1$) approximately
$\theta^{(q)}_k-\theta^{(q)}_{k-1}\sim\Norm\!\big(0,\sigma_\theta^2\,g(\Delta_k)\big)$, with
$g(\Delta)$ the relevant entry of $\bm{Q}_\theta(\Delta)$, so $\sigma_\theta^2$ is recovered by
matching the empirical increment variance (equivalently, by maximising the IWP marginal likelihood of
the solution track). The same procedure on offline flux solutions sets the post-warm-up floor of
$\sigma_A^2$, and the spread of fitted values across observations sets sensible defaults---they change
little between observations, which is precisely why fixing them is safe. The fitter is a small offline
utility emitting a frozen hyperparameter configuration; it never runs inside the streaming pipeline.

\section{Implementation notes}
\label{sec:notes}

The pipeline is realised as a streaming graph of separate operators---reader $\to$ sky-model $\to$
calibration $\to$ imaging $\to$ writer---so each stage has a single responsibility and can be tested
in isolation:
\begin{itemize}[leftmargin=*]
  \item \textbf{Sky-model operator} computes $X_{pq,s}(t)$ from the catalogue and the per-frame
        geometry. TART is single-polarisation, so one scalar coherency per source/baseline suffices.
        The ephemeris/coordinate logic can be borrowed from TART community code (the
        \texttt{tart-telescope} organisation; T.~Molteno's \textsc{disko} for HEALPix imaging and
        \textsc{spotless} for visibility-domain source removal)~\cite{disko}, but those are host/NumPy
        implementations---we re-implement the per-frame kernels as GPU (CuPy) operators rather than
        calling them at runtime.
  \item \textbf{Calibration operator} is stateful (it carries the filter mean/covariance, the gauge,
        and the warm-up schedule across frames), owns the subtraction, and emits the gain-corrected
        residuals \eqref{eq:resid} plus a solutions side-stream (gains, fluxes, NIS) for diagnostics
        and for the offline hyperparameter fit.
  \item \textbf{Imaging operator} is a new HEALPix DFT operator, structured after the existing tiled
        $(l,m)$ operator, consuming residuals in the raw-visibility layout.
  \item \textbf{Reader} is deferred: TART data are not delivered as MSv4 Zarr, and the concrete ingest
        format will be fixed once sample data are in hand. The existing MSv4 reader serves only as a
        structural template.
\end{itemize}
The numerical backend is CuPy throughout, which also lets the calibration and imaging operators
exchange GPU buffers with the existing CuPy/Holoscan tensors at zero copy. Because the calibration
Jacobian is structured (antenna-bipartite gains, dense flux block), the update is hand-coded
structured linear algebra rather than relying on a high-level autodiff framework.

\begin{thebibliography}{99}
\small
\bibitem{kailath} T.~Kailath, ``An innovations approach to least-squares estimation, Part I,'' \emph{IEEE Trans.\ Autom.\ Control}, 13(6):646--655, 1968.
\bibitem{barshalom} Y.~Bar-Shalom, X.-R.~Li and T.~Kirubarajan, \emph{Estimation with Applications to Tracking and Navigation}, Wiley, 2001.
\bibitem{sarkka} S.~S\"arkk\"a and A.~Solin, \emph{Applied Stochastic Differential Equations}, Cambridge University Press, 2019.
\bibitem{stefcal} S.~Salvini and S.~J.~Wijnholds, ``Fast gain calibration in radio astronomy using alternating direction implicit methods,'' \emph{Astron.\ Astrophys.}, 571:A97, 2014.
\bibitem{studentt} M.~Roth, E.~\"Ozkan and F.~Gustafsson, ``A Student's t filter for heavy tailed process and measurement noise,'' \emph{IEEE ICASSP}, 2013.
\bibitem{healpix} K.~M.~G\'orski et al., ``HEALPix: A framework for high-resolution discretization and fast analysis of data distributed on the sphere,'' \emph{Astrophys.\ J.}, 622:759--771, 2005.
\bibitem{tart} The TART collaboration, \emph{Transient Array Radio Telescope}, \url{https://github.com/tart-telescope}.
\bibitem{disko} T.~Molteno, \textsc{DiSkO} and \textsc{spotless}, \url{https://github.com/tmolteno}.
\end{thebibliography}

\end{document}
